% GENERATED FILE. Edit canonical lessons, then run npm run generate:books.
% canonical-source-hash: fnv1a64:67aed8231cfa56c2
% canonical-lessons: IT-C471-timidezza, IT-C471-socievole, IT-C471-riservato, IT-C471-chiacchierone, IT-C471-arrossire

\chapter{Timidi e socievoli}
\label{ch:it-timidi-e-socievoli}
\hlchaptermodality{471}

\begin{chapteropening}
\textbf{By the end of this chapter:} \emph{I can describe how shy or outgoing someone is.}

Complete the last lesson of chapter 471: I can describe how shy or outgoing someone is.
\end{chapteropening}

\section[la timidezza]{la timidezza — shyness}

\label{lesson:IT-C471-timidezza}



\begin{quote}
Before the new one: say the Italian for exactly, then the Italian for almost.
\end{quote}

\subsection*{You'll want to know: la timidezza}
\textbf{la timidezza} — \textquotedblleft{}shyness\textquotedblright{}.

\begin{quote}
Con il tempo ho superato la timidezza.
\end{quote}

\begin{quote}
\emph{With time I got over my shyness.}
\end{quote}

Built on \emph{timido} \textquotedblleft{}shy\textquotedblright{} with \emph{-ezza}, the ending that names a quality, as in \emph{bellezza}.

\begin{cousinweb}
\emph{timidezza} is built on \emph{timido} \textquotedblleft{}shy\textquotedblright{}, a learned borrowing from Latin \emph{timidus} \textquotedblleft{}fearful\textquotedblright{} (Lewis \& Short), from \emph{timēre} \textquotedblleft{}to fear\textquotedblright{}; \emph{-ezza} turns it into the name of a quality.
\end{cousinweb}

\begin{grammarlens}[title={What you've built}]
One more word for telling, explaining and describing at B1.
\end{grammarlens}

\subsection*{Your turn}
\begin{itemize}
\raggedright
  \item \emph{Say it:} \emph{la timidezza}
  \item \emph{Say it:} \emph{la timidezza}, once more
  \item \emph{Say it:} a sentence of your own with \emph{la timidezza}
  \item \emph{Recall:} say the Italian for exactly, then the Italian for almost, then say \emph{la timidezza} again
\end{itemize}

\begin{tcolorbox}[breakable,colback=teal!4,colframe=teal!35!black,title={Before you move on}]
What does la timidezza mean? (\textquotedblleft{}Shyness\textquotedblright{}.) Say it once more.
\end{tcolorbox}
\section[socievole]{socievole — sociable, outgoing}

\label{lesson:IT-C471-socievole}



\begin{quote}
Before the new one: say the Italian for almost, then the Italian for shyness.
\end{quote}

\subsection*{You'll want to know: socievole}
\textbf{socievole} — \textquotedblleft{}sociable, outgoing\textquotedblright{}.

\begin{quote}
Mia sorella è molto socievole: parla con tutti.
\end{quote}

\begin{quote}
\emph{My sister is very sociable: she talks to everyone.}
\end{quote}

The opposite is \emph{poco socievole}; \emph{asociale} is much stronger.

\begin{cousinweb}
\emph{socievole} comes from Latin \emph{sociābilis} \textquotedblleft{}that can be joined together\textquotedblright{} (EtymDB), built on \emph{socius} \textquotedblleft{}a companion, a partner\textquotedblright{} (Lewis \& Short).
\end{cousinweb}

\begin{grammarlens}[title={What you've built}]
One more word for telling, explaining and describing at B1.
\end{grammarlens}

\subsection*{Your turn}
\begin{itemize}
\raggedright
  \item \emph{Say it:} \emph{socievole}
  \item \emph{Say it:} \emph{socievole}, once more
  \item \emph{Say it:} a sentence of your own with \emph{socievole}
  \item \emph{Recall:} say the Italian for almost, then the Italian for shyness, then say \emph{socievole} again
\end{itemize}

\begin{tcolorbox}[breakable,colback=teal!4,colframe=teal!35!black,title={Before you move on}]
What does socievole mean? (\textquotedblleft{}Sociable, outgoing\textquotedblright{}.) Say it once more.
\end{tcolorbox}
\section[riservato]{riservato — reserved, private (of a person); booked}

\label{lesson:IT-C471-riservato}



\begin{quote}
Before the new one: say the Italian for shyness, then the Italian for sociable.
\end{quote}

\subsection*{You'll want to know: riservato}
\textbf{riservato} — \textquotedblleft{}reserved, private (of a person); booked\textquotedblright{}.

\begin{quote}
Mio figlio è un ragazzo riservato e parla poco di sé.
\end{quote}

\begin{quote}
\emph{My son is a reserved boy and says little about himself.}
\end{quote}

Of a table or a seat it means booked: \emph{un tavolo riservato}.

\begin{cousinweb}
\emph{riservato} is the past participle of \emph{riservare} \textquotedblleft{}to keep back\textquotedblright{}, a learned borrowing from Latin \emph{reservāre} (Lewis \& Short), built on \emph{servāre} \textquotedblleft{}to keep, to watch over\textquotedblright{}.
\end{cousinweb}

\begin{grammarlens}[title={What you've built}]
One more word for telling, explaining and describing at B1.
\end{grammarlens}

\subsection*{Your turn}
\begin{itemize}
\raggedright
  \item \emph{Say it:} \emph{riservato}
  \item \emph{Say it:} \emph{riservato}, once more
  \item \emph{Say it:} a sentence of your own with \emph{riservato}
  \item \emph{Recall:} say the Italian for shyness, then the Italian for sociable, then say \emph{riservato} again
\end{itemize}

\begin{tcolorbox}[breakable,colback=teal!4,colframe=teal!35!black,title={Before you move on}]
What does riservato mean? (\textquotedblleft{}Reserved, private (of a person); booked\textquotedblright{}.) Say it once more.
\end{tcolorbox}
\section[il chiacchierone]{il chiacchierone — chatterbox; talkative}

\label{lesson:IT-C471-chiacchierone}



\begin{quote}
Before the new one: say the Italian for sociable, then the Italian for reserved.
\end{quote}

\subsection*{You'll want to know: il chiacchierone}
\textbf{il chiacchierone} — \textquotedblleft{}chatterbox; talkative\textquotedblright{}.

\begin{quote}
Il mio vicino è un gran chiacchierone.
\end{quote}

\begin{quote}
\emph{My neighbour is a real chatterbox.}
\end{quote}

A woman is \emph{la chiacchierona}. It is built on \emph{chiacchierare} \textquotedblleft{}to chat\textquotedblright{}.

\begin{cousinweb}
\emph{chiacchierone} is built on \emph{chiacchierare} \textquotedblleft{}to chat\textquotedblright{} with \emph{-one}, the ending for someone who does a lot of something (EtymDB). The verb itself is usually taken to imitate the sound of chatter.
\end{cousinweb}

\begin{grammarlens}[title={What you've built}]
One more word for telling, explaining and describing at B1.
\end{grammarlens}

\subsection*{Your turn}
\begin{itemize}
\raggedright
  \item \emph{Say it:} \emph{il chiacchierone}
  \item \emph{Say it:} \emph{il chiacchierone}, once more
  \item \emph{Say it:} a sentence of your own with \emph{il chiacchierone}
  \item \emph{Recall:} say the Italian for sociable, then the Italian for reserved, then say \emph{il chiacchierone} again
\end{itemize}

\begin{tcolorbox}[breakable,colback=teal!4,colframe=teal!35!black,title={Before you move on}]
What does il chiacchierone mean? (\textquotedblleft{}Chatterbox; talkative\textquotedblright{}.) Say it once more.
\end{tcolorbox}
\section[arrossire]{arrossire — to blush, to go red}

\label{lesson:IT-C471-arrossire}



\begin{quote}
Before the new one: say the Italian for reserved, then the Italian for chatterbox.
\end{quote}

\subsection*{You'll want to know: arrossire}
\textbf{arrossire} — \textquotedblleft{}to blush, to go red\textquotedblright{}.

\begin{quote}
Quando parlo in pubblico arrossisco subito.
\end{quote}

\begin{quote}
\emph{When I speak in public I go red straight away.}
\end{quote}

An \emph{-isc-} verb: \emph{arrossisco}, \emph{arrossisci}. It takes \emph{essere} in the past: \emph{sono arrossito}.

\begin{cousinweb}
\emph{arrossire} is built from \emph{a-} and \emph{rosso} \textquotedblleft{}red\textquotedblright{}, which is inherited from Latin \emph{russus} \textquotedblleft{}red\textquotedblright{} (Lewis \& Short).
\end{cousinweb}

\begin{grammarlens}[title={What you've built}]
One more word for telling, explaining and describing at B1.
\end{grammarlens}

\subsection*{Your turn}
\begin{itemize}
\raggedright
  \item \emph{Say it:} \emph{arrossire}
  \item \emph{Say it:} \emph{arrossire}, once more
  \item \emph{Say it:} a sentence of your own with \emph{arrossire}
  \item \emph{Recall:} say the Italian for reserved, then the Italian for chatterbox, then say \emph{arrossire} again
\end{itemize}

\begin{tcolorbox}[breakable,colback=teal!4,colframe=teal!35!black,title={Before you move on}]
What does arrossire mean? (\textquotedblleft{}To blush, to go red\textquotedblright{}.) Say it once more.
\end{tcolorbox}
